\scanpage{70}

\deleted{\deleted{(alg clos} variables $k$ \inserted{alg clos (pour simplifier)} de car. première à $n$, \deleted{des} \uline{munis} d'un $\alpha\in k$ tel que les solutions de $\zeta^2-\alpha\zeta+1=0$ soient des racines primitives $n$\textsuperscript{ièmes} de $1$ (ce qu'on peut exprimer par la condition $F_n(\zeta)=0$ \dots) une \deleted{\unsure{rep}} famille de $k$-représentations affines de dim $2$ \inserted{de $D_n$} \unsure{ttes|tous|des} isomorphes -- savoir celles associées aux \deleted{\inserted{représentations} polygônes réguliers} couples formés d'un polygône régulier d'invariant $\alpha$, et d'un épinglage de celui-ci -- ou encore des hom $D_n\longrightarrow\operatorname{Aff}\unsure{(2k)|(2,k)|(3k)}$ définis à autom intérieur près.}

\uline{\unsure{Remarques|Remarque}}

Si on s'était borné d'exiger, pour un polyèdre \inserted{affine} \uline{régulier}\deleted{affine}, que le gpe de ses autom \uline{affines} soit transitif sur l'ens des sommets, \inserted{(sans \unsure{épuiser|préciser|exiger} \unsure{néc|néc\textsuperscript{t}|rien} \unsure{l'ens|leurs} \deleted{\unsure{groupe}} des autom \unsure{comb|combs|ensembles})} on aurait trouvé beaucoup plus de \unsure{classes|choses}. Par exemple, prenant le cas particulier où \inserted{le gpe $G^+$} des rotations combinatoires provient d'un gpe de transformations affines, on aurait eu à classer des $u$ tels que $u^n=1$ (sans disposer de $\sigma$ tel que $\sigma u\sigma^{-1}=u^{-1}$) ce qui nous aurait donné deux val. propres $\zeta,\zeta'$ racines \deleted{\unsure{pri|pr}} $n$\textsuperscript{ièmes} de $1$, d'ordre $d,d'\mid n$, avec

% archive sheet number: 69
